%% ---------- Notation and standing conventions for Cluster D ----------
%% Source: kl_localization_sheet_resolution_factorization.md, Sec. 2-3 (lines 62-183)
\subsection*{Notation and standing conventions}

All spaces in this part are Hausdorff, connected, locally path-connected,
and semilocally simply connected; all covering spaces are connected and
based.  These hypotheses are satisfied by the planar complementary
surfaces used in the canonical KL localization.  Path products and
induced fundamental-group maps are based; every statement involving an
unbased cover or subgroup is to be read after the displayed basepoints
have been fixed.

For a based covering \( p\colon(\widehat X,\widehat x)\to(X,x) \) we write
\[
  H \;=\; p_*\pi_1(\widehat X,\widehat x)\;\leq\;\pi_1(X,x)
\]
for its characteristic subgroup.

For subgroups \( H\leq K \) of a group,
\begin{equation}\label{eq:coset-set}
  K/H=\{\,kH : k\in K\,\}
\end{equation}
denotes only the \emph{pointed set} of cosets with distinguished point
\( H \).  No multiplication on \eqref{eq:coset-set} is asserted unless
\( H\triangleleft K \).  With the usual path-endpoint convention a cover
fiber is naturally the right-coset set \( H\backslash K \); inversion
gives the canonical pointed bijection \( Hk\longmapsto k^{-1}H \) with
\eqref{eq:coset-set}.  Thus the notation \( K/H \) does not silently
assume normality.

For composable paths we use the paper-wide convention of
Section~\ref{sec:background}: \( \delta_1*\delta_2 \) means first
traverse \( \delta_1 \), then traverse \( \delta_2 \), and in
fundamental groups \( [\delta_1][\delta_2]=[\delta_1*\delta_2] \)
(the first-traversed factor is written leftmost).  In the endpoint
description, continuing the fiber point labelled \( H_sk \) by a based
loop \( g \) appends \( g \) after \( k \) and gives \( H_skg \); under
the inversion identification
\( H_sk\mapsto k^{-1}H_s \), that transport is
\begin{equation}\label{eq:transport-formula}
  k^{-1}H_s\;\longmapsto\; g^{-1}k^{-1}H_s .
\end{equation}
Formula \eqref{eq:transport-formula} replaces any naive multiplication on
\( K/H_s \) and remains valid when \( H_s \) is not normal in \( K \).

Fixed symbols of this part:
\( a\colon(X_s,x_s)\hookrightarrow(X_t,x_t) \) is a based topological
embedding (at the canonical KL primitive seed the spaces are the
complementary surfaces, \( X_s=\Omega_s \), \( X_t=\Omega_t \), and
\( a=i^{p} \));
\( p_s\colon(\widehat X_s,\widehat x_s)\to(X_s,x_s) \) and
\( p_t\colon(\widehat X_t,\widehat x_t)\to(X_t,x_t) \) are connected based
covers with subgroups \( H_s,H_t \);
\( A\colon(\widehat X_s,\widehat x_s)\to(\widehat X_t,\widehat x_t) \) is a
based lift satisfying \( p_tA=ap_s \);
\( K:=a_*^{-1}(H_t)\leq\pi_1(X_s,x_s) \);
\( P=X_s\times_{X_t}\widehat X_t \) is the pullback and \( P_K \) the path
component of the basepoint \( (x_s,\widehat x_t) \), with first and second
projections \( r \) and \( e \);
\( q\colon\widehat X_s\to P_K \) is the canonical comparison map of
Theorem~\ref{thm:sheet-resolution};
\( \mathcal S_{\mathrm{loc}}=K/H_s \) is the based localization-sheet
fiber;
\( M_\delta \) denotes sheet transport along a path \( \delta \) in
\( P_K \); \( \Pi_1(P_K) \) is the fundamental groupoid of \( P_K \); and
\( \mathcal G_q \) is the canonical sheet-transport groupoid of
Proposition~\ref{prop:transport}.  \( \mathbf{Set}_{\mathrm{bij}} \)
denotes the category of sets and bijections.  A map defined on an open
interval is \emph{topologically proper} if preimages of compact sets are
compact.

%% Typographic note: the source note writes \mathcal{S}_{loc} and
%% \mathcal{G}_q; we use \mathcal{S}_{\mathrm{loc}} and \mathcal{G}_q to
%% remain within amssymb (no mathrsfs).  The source's amscd square in
%% Corollary 7.2 is rendered as a displayed commutation equation.